\otherchapter{观测空间与策略网络设计}
\section{传感器配置与状态表示}
\label{sec:sensor_state}

本文策略网络的输入信息完全依赖于机器人自身的低延迟传感器。我们将所有必要的感知数据组合成一个状态特征向量 $\mathbf{s}_t$，其构成可表述如下：

\begin{equation}
\mathbf{s}_t = [\mathbf{g}, \mathbf{v}^d, \bm{\phi}, \dot{\bm{\phi}}, \mathbf{a}_{t-1}] \in \mathbb{R}^{n_s}
\label{eq:state_vector}
\end{equation}

这个向量包含了让机器人了解自身处境和任务目标的关键信息。其中，$\mathbf{g} \in \mathbb{R}^3$ 是由机载惯性测量单元（IMU）提供的重力方向向量，它直接反映了躯干相对于重力场的姿态，是判断平衡状态的核心依据，相比使用完整的姿态四元数，这种方式显著降低了计算开销。$\mathbf{v}^d = [v_x^d, v_y^d, \omega_z^d]$ 是高层系统下达的速度指令，指明了期望的前进速度、侧向速度和转向速度，将指令融入状态使得单一策略网络能够适应多种运动任务，避免了为每个速度单独训练网络的需要。

$\bm{\phi}$ 和 $\dot{\bm{\phi}}$ 分别代表了所有关节的当前位置和角速度，它们定义了机器人的实时几何构型与运动趋势，是做出任何动力学合理决策的基础。对于本研究涉及的12自由度模型，关节相关状态为24维；对于29自由度全身模型，则为58维。此外，状态中还包含了上一时刻的动作输出 $\mathbf{a}_{t-1}$。引入这一“历史记忆”主要是为了鼓励动作在时间上的平滑性，使网络在决策时能参考先前的动作，从而生成连贯的运动序列，有效抑制各时间步独立决策可能引发的高频抖动。

综合来看，策略网络所接收的完整状态向量，其维度由前述的五个部分共同决定。具体而言，它包含了反映机体姿态的3维重力向量、3维基座角速度、以及分别与关节数量对应的位置、速度与历史动作信息。因此，对于只控制下肢的12自由度简化模型，其状态向量的总维度便是各部分维度之和，即3维重力、3维角速度、12维关节位置、12维关节速度与12维历史动作，共计42维。相应地，在需要协调全身运动的29自由度模型中，由于驱动关节数目增加，关节位置、速度与历史动作这三个部分的维度均提升至29，使得状态向量的总维度扩展为93维。这个包含了机器人瞬时状态与近期动作历史的高维向量，随后被输入到多层感知机构成的策略网络中，网络最终输出的动作指令维度，则严格对应于当前模型需要控制的可动关节总数。

为了提高训练的稳定性与效率，我们对状态向量进行了标准化预处理。具体采用了一种“冻结统计量”的方法：在训练初期采集一小批样本数据，计算各状态维度的均值 $\bm{\mu}$ 和标准差 $\bm{\sigma}$，并在整个训练过程中固定使用这些统计量对输入状态进行标准化，即 $\mathbf{s}_t^{\text{norm}} = (\mathbf{s}_t - \bm{\mu}) / \bm{\sigma}$。这种做法能够有效防止训练过程中因数据分布漂移而引发的不稳定，并将不同物理意义和量纲的状态统一到相近的数值尺度，从而确保梯度优化算法能够平稳、高效地工作。如果缺少这一步，网络很容易对某些数值范围较大的状态分量过度敏感，导致训练收敛缓慢或最终策略表现不佳。因此，输入标准化是保障深度强化学习算法有效利用网络容量、成功学习策略的关键步骤。


\section{动作映射与执行控制}
网络输出的动作集合并非直接对应关节力矩，而是经缩放操作后叠加于基础站立参考位姿。采用相对坐标设计的优势在于：首先，直接让网络学习绝对关节角度会使策略偏袒处于特定状态的绝对位置而轻视系统动态响应；其次，初始化差异会促使网络收敛落差变大，削弱了策略泛化性能。基于“自然站立位姿”的基础微调，能够促使策略网络具备快速趋于稳定动作的能力。这一学习范式有助网络将优化焦点汇聚于动作的“增量需求”上，而非锁定一个精确的预期位置。

具体形式为：
\begin{equation}
q_{ref,t} = q_{default} + k_{scale} \cdot a_t
\label{eq:action_map}
\end{equation}
其中 $q_{default}$ 是人为定义的标准站立姿态（通常是平衡状态下的关节角度参考）。在代码实现中，$k_{scale}$ 对应配置项 `action_scale`（示例值 `action_scale = 0.25`）；不同模型/实验可在配置中调整该缩放因子。经过这个映射的关节参考角度 $q_{ref,t}$ 再被输入底层的 PD 反馈控制器：
\begin{equation}
\tau_t = K_p (q_{ref,t} - q_{actual,t}) - K_d \dot{q}_{actual,t}
\label{eq:pd_control}
\end{equation}
其中 $K_p$ 和 $K_d$ 为比例和微分增益系数。经大批量并发仿真调优，最终确定所有关节使用统一增益参数：$K_p = 30 \text{ N·m/rad}$，$K_d = 2 \text{ N·m·s/rad}$。过小的增益系数将削弱控制器带宽并引发控制延迟问题，致其无法敏捷追踪网络动作的决策指令；而过高的力矩增益则容易激起系统的高频震荡。在拥有 1024 个并行环境的 Isaac Gym 中进行反复测试验证后，该参数被证实具有较好的控制鲁棒性，能够在保持追踪响应速率的同时规避执行震荡及系统不平稳。
\end{equation}
其中 $K_p$ 和 $K_d$ 为比例和微分增益系数。实际实现中并未使用单一的全局标量 $K_p,K_d$，而是在配置文件中对不同关节分别设置刚度（`stiffness`）和阻尼（`damping`）字典以获得更细粒度的跟踪性能。例如在仓库配置 `legged_gym/envs/g1/g1_config.py` 中，典型值为：髋/膝关节刚度较大（如 hip_pitch≈100 N·m/rad，knee≈150 N·m/rad），踝关节较小（ankle≈40 N·m/rad）；阻尼项在 1–4 范围内（例如 hip_yaw≈2，knee≈4，ankle≈2）。这种按关节细化的做法既能提高跟踪带宽又能防止局部高频震荡。具体实验时使用的配置文件路径与具体数值已记录于代码仓库，便于复现。

这个两层架构（上层 DRL 决策 + 下层 PD 执行）的设计有三个实际优势。首先是解耦——强化学习负责高层决策逻辑，而人工调参的 PD 控制器负责低层跟踪精度，两者各司其职。其次是鲁棒性增强。即使策略在仿真中学到的动作不够精确，下层 PD 控制器仍能通过反馈环节来补偿误差。这在实际部署中特别重要，因为硬件和仿真总会有偏差。第三是便于域转移。从仿真迁移到真实机器人时，主要调整 PD 参数（在真实硬件上很快），而网络权重无需重新训练，大幅加快了部署周期。

本研究实现的策略网络 $\pi_\theta$ 采用基于循环-演员-评论家（recurrent actor-critic）的结构以处理时序依赖：策略头包含 LSTM 层用于聚合历史观测，随后接小型前馈网络负责最终的动作参数化。实际训练配置（见 `legged_gym/envs/g1/g1_config.py`）使用的网络尺寸为：LSTM 隐藏维度 `rnn_hidden_size = 64`，前馈隐藏层为 `actor_hidden_dims = [32]` 与 `critic_hidden_dims = [32]`。网络输出仍以高斯分布参数化（均值与对数标准差），动作通过采样获得。该配置是在大规模并行训练（例如 1024 个环境）下，在计算效率与时序建模能力之间取得的折中；如需更大模型（例如隐藏 256 维），请在配置中调整相应字段并在实验中重新评估训练稳定性与耗时。

\section{奖励函数设计}
我们设计了多目标奖励函数 $\mathcal{R}_{total}=\sum_{i}w_{i}r_{i}$，需要平衡多个指标。单纯追求速度会导致"原地快踏步"这样的无效动作。只限制平滑度会使步态拖沓。所以必须通过权重和训练流程的调整来平衡这些目标。速度跟踪：使用高斯奖励函数
\begin{equation}
r_{track\_v}=\exp\left(-\frac{||v_{actual}-v^{cmd}||^{2}}{2\sigma_{v}^2}\right), \quad \sigma_v = 0.5
\label{eq:velocity_reward}
\end{equation}
奖励函数的设计需要引导机器人学习到既高效又稳定的步态。其中，位置跟踪奖励 $r_{\text{vel}}$ 采用了Huber损失函数的形式，这样的设计是考虑到，当跟踪误差比较大的时候，损失函数的梯度会受到限制，从而能够防止训练过程因为个别异常的大误差样本而产生不稳定；而当跟踪误差很小时，它仍然能提供一个有效的梯度，来驱动策略网络继续优化、朝着更精确的方向收敛。对于机器人的转向控制，角速度跟踪奖励 $r_{\text{ang\_vel}}$ 也采用了类似的形式，以稳定机器人的整体朝向。为了引导机器人进行更节能的运动，我们引入了力矩惩罚项 $r_{\text{torque}} = -\sum_i \tau_{i}^2$。这项惩罚直接对各个关节输出力矩的平方和进行负向激励，本质上是鼓励策略网络在满足运动要求的前提下，尽可能地去减少整体的能量消耗和电机发热。同时，为了避免策略输出那些高频抖动的、不自然的动作序列，我们特别设置了动作平滑性惩罚 $r_{\text{act\_smooth}} = -||\mathbf{a}_{t} - \mathbf{a}_{t-1}||^{2}$，它会对相邻两个时间步的动作向量之间的剧烈变化进行惩罚。我们在实验中观察到一个有趣的现象：当这项惩罚的权重被设置为 $-0.01$ 时，机器人步态的整体抖动幅度能够减少大约 $35\%$，得到的跨步动作看起来也明显更连贯、更接近真实的步态关节的运动范围是硬件上必须遵守的硬性约束。我们通过关节限幅惩罚 $r_{\text{limit}}$ 来实现这一点：只要检测到任何一个关节的角度 $q_i$ 超过了其物理限位 $q_{i, \text{max}}$，就会立即施加一个较强的固定惩罚（实验中设为 $r_{\text{limit}} = -1.0$）。这相当于一个“红线”警告，迫使策略网络必须在学习过程中主动地“学会”并严格遵守这些机械边界此外，为了塑造出更规律、更稳定的基本步态模式，我们还设定了针对躯干姿态和足底高度的约束规则。当机器人的躯干倾斜角度超过30度时，会触发一个姿态惩罚，这是为了防止它发生严重的倾倒。同时，如果机器人的脚在摆动阶段离地时间过长，也会受到相应的惩罚，这有助于引导它形成周期分明、触地时间合理的步态整个奖励函数的调试过程本身也给我们带来了一些启发。我们发现，如果在训练早期只使用速度跟踪奖励，机器人会很快学会一种步频极高、但移动效率很低的“快踏步”，这更像是一种针对奖励规则的“投机”行为。而及时地引入动作平滑惩罚后，步态质量得到了显著的改善。另一方面，像关节限幅、姿态约束这类关键的惩罚项，如果引入得太晚，策略网络可能已经固化了一些不良的动作习惯，到后期再想纠正就会变得非常缓慢。因此，各项奖励和惩罚的权重，需要与核心的速度奖励进行协同调整，最终目的是在“鼓励机器人前进”与“约束其动作质量”之间，找到一个好的平衡点下表总结了本文最终采用的多维度复合奖励函数的具体构成与经过调试后的权重配比。

\begin{table}[h]
\centering
\caption{核心奖励函数设计与权重配比}
\begin{tabular}{llcc}
\toprule
\textbf{奖励类别} & \textbf{设计意图} & \textbf{权重 w} & \textbf{物理涵义} \\
\midrule
\multirow{2}{*}{速度跟踪} & 线速度跟踪 & +2.0 & 促使机器人前行或侧移符合指令 \\
         & 角速度跟踪 & +0.5 & 保障对转向指令的快速响应 \\
\midrule
\multirow{1}{*}{姿态与轨迹约束} & 维持基础身高 & -10.0 & 惩罚质心过低或下蹲行走 \\
         & 约束脚摆高度 & -20.0 & 规范跨步时的抬脚离地轨迹 \\
\midrule
\multirow{1}{*}{能耗与平滑} & 关节加速度平滑 & -2.5e-7 & 抑制运动中出现的异常加速度抖动 \\
\bottomrule
\end{tabular}
\end{table}

奖励函数中各部分的权重设置，并不是单纯通过理论推导一次性确定的，而是在Isaac Gym仿真环境中，通过反复观察机器人的实际训练表现，经过多次调试才最终确定的。在训练的初期阶段，我们注意到，如果没有对关节的输出力矩加以有效限制，机器人会为了尽可能地获得“速度跟踪”这一项的奖励，逐渐演化出一种步态异常——它的步频变得极高，步幅却非常小，腿部几乎是在地面上快速地“拖行”或“蹭动”（这种现象在足式机器人中常被称为Scuffing）。这显然不是我们期望看到的、能量利用效率高且符合物理规律的行走方式。

为了解决这个问题，我们对奖励函数进行了针对性的调整。一方面，我们显著提高了动作平滑性奖励项 $r_{\text{act\_smooth}}$ 的权重（设置为$-0.01$），以惩罚关节角度在相邻时间步之间的剧烈变化，鼓励生成更流畅的动作序列。另一方面，我们同时强化了关节扭矩惩罚项 $r_{\text{torque}}$ 的作用。这两项奖励的共同引导，迫使策略网络在寻求高移动速度的同时，必须综合考虑动作的连贯性与能量消耗，从而最终找到了那些能量利用率更高、也更符合真实关节电机物理特性与工作限值的跨步轨迹。

